\section{Exit paths, torus links, and the exact criterion}
\label{sec:criterion}

\subsection{The geometric datum behind the groupoids}

\begin{definition}[\status{Proved}: the exit category and its link model]\label{P1a1}
For a normal-crossings pair in the preceding scope, $\Exit(X,D)$ is the
exit-path $\infty$-category of the stratification by connected open strata.
On a manifold it is the stratified singular simplicial set: the last-vertex
stratification on each simplex is required to map increasingly to the
frontier order.  On a stack it is the colimit of these categories on an
étale presentation.  Its restriction to a stratum $S$ is $\Pi_\infty S$.
For incident strata $S<T$, of complex codimensions $d>e$, let $F_{ST}\to S$
be the finite cover that chooses the local branch of $\overline T$ along
$S$.  There is a torus bundle
\[
 (S^1)^{d-e}\longrightarrow L_{ST}\xrightarrow{p}F_{ST},
 \qquad (a,b):L_{ST}\longrightarrow S\times T,                 \tag{1}
\]
where $a$ is the projection and $b$ is a small positive radial displacement,
well defined up to a coherent homotopy.  For geometric points $s\in S$,
$t\in T$,
\[
 \Map_{\Exit(X,D)}(s,t)\simeq
 \hofib_{(s,t)}\bigl(L_{ST}\xrightarrow{(a,b)}S\times T\bigr). \tag{2}
The right side is the disjoint union over \emph{all} incidence components
and endpoint transports.  For $S=T$ it is the ordinary path space in $S$.
The bundle in (1), not the fixed-endpoint mapping space in (2), is the
geometric object represented by $\breve N^{**}$ in 147/177--178, (295).
\end{definition}
\begin{proof}
Near a point with branch set $I$, use normal-crossing coordinates
$(z_i)_{i\in I},w$, with $D=\{\prod z_i=0\}$.  A local branch of $T$ is a
subset $J\subset I$, $|J|=e$, of coordinates that stay zero.  An immediate
exit to $T$ has $z_j=0$ for $j\in J$ and $z_i\ne0$ for $i\in I\setminus J$.
Its normal-direction space has a deformation retraction onto
$(S^1)^{I\setminus J}$: first divide by the positive radial coordinate, then
contract the positive ratios of the $d-e$ absolute values.  The latter
parameter space is the interior of a simplex and is contractible.  Changes
of adapted coordinates permute the normal lines and multiply them by
nonzero transition functions.  Thus these local direction spaces glue to
exactly the product of their unit normal line bundles over $F_{ST}$.

We explain why using immediate exits loses no mapping-space information.
The normal-crossing stratification is a Whitney stratification and has
control data obtained by the coordinate radial contractions.  These give
neighbourhood tameness and the fibration of the homotopy-link evaluation
over the source stratum.  Miller's endpoint-preserving comparison
\citep[Theorem 4.9]{Miller2009} deforms the popath space to the homotopy
link; the theorem applies to the induced finite interval of strata between
$S$ and $T$.  The proof is local over the source and is compatible with
restriction, so locally finite pairs are covered as well.  Equivalently,
successively shorten the intervals spent in intermediate coordinate
strata and use the radial contraction to make all coordinates in
$I\setminus J$ nonzero immediately.  The parameter spaces of positive
shortening functions and radii are contractible.  This is a weak
equivalence over both endpoint spaces.  Replacing endpoint evaluation by
a fibration gives (2), rather than an incorrect ordinary fibre of a
chosen, non-fibrant link map.

For a stack perform these constructions on every term of its étale Čech
nerve.  Branch choices, normal line bundles, and endpoint evaluations all
pull back under étale maps.  Descent of constructible sheaves identifies
the resulting colimit with the stack exit category
\citep[Proposition 3.16 and Corollary 3.17]{JansenToolbox2024}; applying this also to the
endpoint diagrams gives the stack version of (2).  Finite stabilizers
are not discarded, and a finite cover $F_{ST}\to S$ is interpreted as a
representable covering of stacks.  The manifold exodromy and
groupoidification assertions are
\citet[Theorem A.9.3 and Corollary A.9.4]{LurieHA2017}: their hypotheses
hold because normal-crossing coordinate charts are conical, locally
contractible, and have finite depth.  This also fixes the direction of
the exit arrows.
\end{proof}

\subsection{What the Remarque actually reconstructs}

\begin{theorem}[\status{Proved}: three distinct reconstruction statements]\label{P1c1}
Let $\mathsf T$ be the \emph{whole coherent} diagram of elementary pieces
of 147/163--178, including all flags, endpoint maps, and their composition
coherences.  For a normal-crossings manifold or stack:
\begin{enumerate}[label=(\alph*)]
\item Without an asphericity or meridian-injectivity hypothesis,
 $\Pi_1\mathsf T$ determines the fundamental groupoids of $X$, its
 locally closed unions of strata, and their tubes and punctured tubes.
 It also determines their ordinary categories of constructible sheaves
 of sets.  In particular it determines the quasi-Galois topos meant in
 147/179--180.
\item If all strata are aspherical, the same coherent diagram determines
 the \emph{actual homotopy types} of these spaces by homotopy colimits
 of the classifying spaces of its groupoids.  No injectivity assumption
 is needed for this assertion.
\item If all strata are aspherical and (6) holds, it furthermore
 determines their constructible $\infty$-topoi as functor categories on
 ordinary exit categories; the shapes of these $\infty$-topoi and of
 their associated ordinary presheaf topoi agree with the actual
 homotopy types.  Thus every homotopy assertion of the Remarque is
 valid in this precise interpretation under the criterion.
\end{enumerate}
The word ``hence'' in the passage from the ordinary constructible topos
to the actual homotopy type is false under asphericity alone.  Conversely,
the claim that meridian injectivity is \emph{necessary} for all derived
reconstruction from the groupoid system is false.  A strict colimit of
groups, a bicategorical colimit of groupoids, a homotopy colimit of their
classifying spaces, and an ordinary constructible topos are different
constructions.
\end{theorem}
\begin{proof}
For two strata $S$ closed and $T$ open, the exit datum consists of the
two fundamental $\infty$-groupoids and the crossing bimodule in (2).
A constructible sheaf of sets consists of local systems $F_S,F_T$ and
a map $a^*F_S\to b^*F_T$ on the link.  Local systems of sets on any
locally contractible space or étale stack are functors from its
fundamental groupoid.  The same is true of their natural
transformations on the link.  These data and their descent condition
therefore use only $\Pi_1 S$, $\Pi_1 T$, and $\Pi_1 L_{ST}$ with
their maps.  With more strata, the equality of two composites is
imposed on each flag space; all the necessary data are included in
$\Pi_1\mathsf T$.  This reconstructs $h\Exit$ and
\[
 \operatorname{Shv}_{\rm cbl}(X;\mathbf{Set})
      \simeq\operatorname{Fun}(h\Exit(X,D),\mathbf{Set}).    \tag{12}
\]
The same argument on a locally closed union or elementary tube gives
the stated canonical topoi.  No assertion about the whole topological
topos of \emph{all} sheaves, or about its ringed analytic structure,
is being made.

The homotopy-colimit decomposition of the next theorem gives the
fundamental groupoid assertion.  Indeed $\tau_{\leq1}$ is left
adjoint to the inclusion of 1-types, so it preserves colimits in
1-types.  Thus $\Pi_1\hocolim\mathsf T$ is the 2-colimit of
$\Pi_1\mathsf T$.  This includes the transport isomorphisms and
the flag relations; a colimit of chosen one-object groups would
discard some of them.

If the strata are aspherical, every elementary piece is a torus
bundle over a finite cover of a stratum, with some contractible disc
factors.  Its homotopy exact sequence makes it a 1-type.  The natural
map $\mathsf T\to B\Pi_1\mathsf T$ is then an objectwise equivalence
of coherent diagrams, so its homotopy colimit is an equivalence.
This proves (b).  There is no replacement of a homotopy fibre by a
set at this step.

Under (6), Theorem~\ref{P1b1} gives
$\Exit(X,D)\simeq N h\Exit(X,D)$, and the same calculation applies
to its restrictions, links, and flag pieces.  Exodromy gives
$\operatorname{Shv}_{\rm cbl}(X;\mathcal S)\simeq
\operatorname{Fun}(h\Exit,\mathcal S)$; the stack case follows from
\citet[Proposition 3.24]{JansenToolbox2024}.  The shape of a functor
$\infty$-topos on an ordinary category $C$ is $BC$, since for every
space $K$ its global sections of the constant object are
$\lim_C K=\operatorname{Map}(BC,K)$.  Groupoidification of exit paths
is $X$ by the theorem cited in the proof of~\ref{P1a1}.  This proves
(c).  Examples~\ref{P1c2} and~\ref{P1d2} prove both negative
assertions and explain precisely which inference the Remords rejects.
\end{proof}

\begin{example}[\status{Refuted}: dropping either hypothesis and reading the margin]\label{P1c2}
There are independent failures of the two hypotheses, with different
consequences for the Remarque (147/179, fol.~85, and 147/180).
\end{example}
\begin{proof}
For $(\PP^1,\{\infty\})$, both strata are contractible, but its
meridian $\Z\to\pi_1(\C)=1$ is zero.  The coherent elementary
diagram has homotopy types $*\leftarrow S^1\to *$, so the
homotopy colimit of its fundamental-groupoid classifying spaces is
$\Sigma S^1\simeq S^2$.  Its ordinary exit category is $[1]$,
whose classifying space is contractible.  Thus only condition (ii)
fails; the ordinary constructible topos and the fundamental groupoid
remain reconstructible, and even the actual homotopy type remains
reconstructible by (b).  It is its identification with the shape of
that ordinary topos that fails.

For $X=\PP^1\times\C$, $D=\PP^1\times\{0\}$, condition (ii)
holds: the meridian maps isomorphically to the $\Z$ factor of
$\pi_1(U)=\Z$.  Condition (i) fails because both $D$ and $U$
have $\pi_2=\Z$.  Its elementary fundamental-groupoid diagram is
identical to that for $(\C^2,\C\times\{0\})$, including the
codimensions and all maps: the base group is trivial and the open
and link groups are $\Z$.  Their ambient homotopy types are $S^2$
and a point.  Thus no reconstruction from that diagram can recover
the ambient homotopy type in the absence of (i).  Their ordinary
constructible topoi and fundamental groupoids are nonetheless
correctly recovered by (a).

The marginal phrase ``indépendt de tt autre hypothèse'' can therefore
be valid for the fundamental-groupoid and constructible-set claims,
not for a claim equating all higher homotopy types with ordinary
presheaf-topos shapes.  We retain the word ``quasi-galoisien'' visible
on scan 179 and ``1-descriptif'' on scan 180; the mathematics does not
justify silently deleting either qualification.  The homology of
the suspension model and of $N[1]$ is independently checked in
\texttt{criterion\_checks.py}.
\end{proof}

\subsection{Gluing, and the homotopy lost by ordinary exit paths}

\begin{theorem}[\status{Proved}: the answer to the gluing question]\label{P1d1}
If a normal-crossings stratification has one closed stratum $S$ and
open complement $T$, its homotopy type and its constructible
$\infty$-topos have shape
\[
 X\simeq\hocolim\bigl(S\xleftarrow{a}L_{ST}
                               \xrightarrow{b}T\bigr).     \tag{13}
\]
For finitely many strata, iterate this formula, retaining crossing
categories at each step, or equivalently realize the full flag
diagram of strata and their iterated homotopy links.  The same formula
holds for the stack homotopy type by descent.  In contrast, the
ordinary constructible-set topos has shape $B h\Exit(X,D)$, obtained
by first taking connected components of the fixed-endpoint mapping
spaces, not by merely replacing $L_{ST}$ by $B\Pi_1L_{ST}$.
\end{theorem}
\begin{proof}
Choose an open tube $N$ of $S$ that retracts onto $S$.  Its intersection
with $T$ is a punctured tube and has type $L_{ST}$ by
Proposition~\ref{P1a2}.  The open cover $X=N\cup T$ has homotopy
Čech realization equivalent to $X$: a partition of unity subordinate
to this two-element cover gives a map into its double mapping
cylinder and a homotopy inverse.  This proves (13), with the maps
specified by the tube projection and the inclusion into $T$.
The same argument can be made in the tubular germ, so different
admissible tubes give the same coherent cocone.

For a finite filtration, apply this two-piece construction to its
successive closed and open parts.  An iterated crossing consists of
a chain of strata and compatible exits; hence all multiple
intersections are exactly the flag spaces.  Associativity of
homotopy colimits gives the flag realization.  Locally finite
stratifications are the filtered union of these constructions on
relatively compact chart unions, and the comparison is checked on
maps from finite spheres and discs, whose images are compact.
For a stack, apply the argument to the étale hypercover and realize;
iterated colimits commute.  This proves the assertion in the stated
generality without assuming a globally embedded stack tube.

Exodromy identifies the constructible $\infty$-topos with
$\operatorname{Fun}(\Exit,\mathcal S)$, whose shape is $B\Exit=X$,
by the constant-object computation in the preceding proof.
Replacing its indexing category by $h\Exit$ gives the distinct
ordinary-presheaf shape.  In the two-stratum case with aspherical
pieces, write
$G=\pi_1(L_{ST})$, $A=\pi_1(S)$, $B=\pi_1(T)$ on one incidence
component.  The original crossing space is $BG$, whereas the
crossing-arrow groupoid for the ordinary category has type
$B\operatorname{im}\{G\to A\times B\}$: its transitive action on
the set of mapping components has exactly this stabilizer.
The quotient by the kernel in (8) is the precise additional step.
For nonconnected incidences use their disjoint union.  This also
derives the appropriate gluing assertion of STRATA C.1--C.3 in the
form needed here, rather than importing its unaudited notation.
\end{proof}

\begin{example}[\status{Proved}: the three requested homotopy fibres]\label{P1d2}
Let $q:X\simeq B\Exit(X,D)\to B h\Exit(X,D)$ be the truncation
comparison.  The following are homotopy fibres over the unique
component of the target:
\[
\begin{array}{c|c|c|c}
(X,D)&B h\Exit(X,D)&\hofib(q)&\text{first failed group}\\ \hline
(\PP^1,\infty)& *&S^2&\pi_2=\Z\\
(\PP^n,H),\ n\geq1& *&\PP^n&\pi_2=\Z\\
((\PP^1)^2,\{\infty\}\times\PP^1\cup\PP^1\times\{\infty\})
 & *&S^2\times S^2&\pi_2=\Z^2
\end{array}                                                   \tag{14}
\]
Only (ii) fails in the first and third rows.  In the second row it
fails for every $n$, and (i) also fails when $n\geq2$.
\end{example}
\begin{proof}
In the first row take objects $p,u$.  Their endomorphism spaces are
contractible and $\Map(p,u)=\hofib(S^1\to\C)=S^1$; all arrows
in the reverse direction are absent.  Thus $h\Exit=[1]$ and its
classifying space is contractible, while $X=S^2$.

For the second row the two strata are $H=\PP^{n-1}$ and
$U=\C^n$.  The crossing total space is the unit normal circle
bundle of $\mathcal O_H(1)$, namely $S^{2n-1}$, up to orientation
of the Hopf fibre.  Equation (2) gives a circle as its fixed-endpoint
mapping space, since $U$ is contractible and the Hopf projection
has fibre $S^1$.  The endomorphism space on the closed stratum is
$\Omega\PP^{n-1}$ and on the open stratum is contractible.  Both
strata are simply connected, so all automorphism \emph{groups} in
$h\Exit$ are trivial and its only cross-stratum Hom-set is a
singleton.  Again $h\Exit=[1]$.  The fibre of a map to a contractible
space is its source, and the standard affine-cell filtration of
$\PP^n$ gives $\pi_2=H_2=\Z$.

Exit paths and their homotopies in a product are coordinatewise,
so the third exit category is the product of two copies of the
first.  Its four strata are $\C^2$, two copies of $\C$, and a
point.  A mapping space that exits in $j$ factors is $(S^1)^j$;
all same-stratum endomorphism spaces are contractible.  Its ordinary
category is $[1]\times[1]$, with terminal object $(u,u)$.
All meridians map to the trivial group of $\C^2$.
The fibre is therefore $(S^2)^2$, and the product CW complex has
two independent two-cells.  These finite chain complexes and all
displayed ranks are checked by the accompanying program.
\end{proof}

\begin{theorem}[\status{Proved}: a general first-obstruction calculation]\label{P1d3}
For any small $\infty$-category $\mathcal C$, the canonical map
$q:B\mathcal C\to B h\mathcal C$ is 2-connected: it is bijective
on $\pi_0$, isomorphic on $\pi_1$, and surjective on $\pi_2$.
Its homotopy fibre is simply connected.  More generally, suppose
every nonempty component of every mapping space is $(r-1)$-connected,
$r\geq1$.  Then $q$ is $(r+1)$-connected.

Here is an explicit first-obstruction module, which also specifies
what ``the first group lost'' can mean without assuming a false
universal numerical degree.  Choose a Kan-enriched category
representing $\mathcal C$, with discrete object set, and put
\[
 V_p=\coprod_{x_0,\ldots,x_p}\prod_{i=1}^p
         \Map_{\mathcal C}(x_{i-1},x_i),\qquad
 W_p=\pi_0V_p=N_p(h\mathcal C).
\]
On a component of $Y=|W_\bullet|$, pull back its universal cover
to the realizations, and denote the resulting degreewise spaces
by $\widetilde V_p,\widetilde W_p$.  Equivalently, label each
component of $V_p$ by the lifts of its simplex in $Y$.
If $F=\hofib(q)$, then
\[
 \pi_{r+1}(F)\cong
 \operatorname{coker}\left[
 H_r(\widetilde V_2;\Z)
 \xrightarrow{d_0-d_1+d_2}
 H_r(\widetilde V_1;\Z)\right].                            \tag{15}
\]
The face $d_1$ is composition; the other two faces are endpoint
projections, with the indicated deck transports.  This is a module
over $\pi_1Y$.  The exact sequence
\[
 \pi_{r+2}(X)\to\pi_{r+2}(Y)\to\pi_{r+1}(F)
 \to\pi_{r+1}(X)\to\pi_{r+1}(Y)\to0                      \tag{16}
\]
identifies the first possible failure, where $X=B\mathcal C$.
For a normal-crossing pair with aspherical strata, $r=1$ and the
generators in (15) are precisely the failed meridian lattices in
(8), transported and subjected to the flag composition relations.
Without additional hypotheses the first failed homotopy group
can occur in any prescribed degree $m\geq2$.
\end{theorem}
\begin{proof}
A Kan-enriched presentation exists by
\citet[Theorem 2.2.5.1]{LurieHTT2009}.  Replacing it by a cofibrant,
well-pointed presentation does not change either realization.
The map $V_p\to W_p$ is $(r)$-connected for $p>0$, since its
fibres are $(r-1)$-connected; it is an isomorphism for $p=0$.
The skeletal filtration of realization shifts the relative cells
in degree $p$ upward by $p$.  Relative cells for these degreewise
maps start in dimension $r+1$, so the relative cells of $q$ start
in dimension $r+2$.  This proves $(r+1)$-connectivity.  This
cellular argument can equally be applied to singular bisimplicial
sets, where latching maps are monomorphisms; it therefore requires
no CW hypothesis on the original mapping spaces.  Taking $r=1$
proves the unconditional assertion.

The same filtration on the mapping-cone chain complex in the
universal covers gives the relative homology spectral sequence
\[
 E^1_{p,j}=H_j(\widetilde W_p,\widetilde V_p;\Z)
       \Longrightarrow H_{p+j}(\widetilde Y,\widetilde X;\Z).
\]
For $p=0$ every group is zero.  For $p>0$ the groups with
$j\leq r$ vanish, and the first row is
$E^1_{p,r+1}=H_r(\widetilde V_p;\Z)$, because
$\widetilde W_p$ is discrete.  At total degree $r+2$ only the
entry $(p,j)=(1,r+1)$ occurs.  Its sole incoming differential is
the displayed alternating sum from $p=2$; all higher incoming
differentials start in the zero rows.  There are no outgoing
differentials.  Consequently this cokernel is
$H_{r+2}(\widetilde Y,\widetilde X;\Z)$.
Relative Hurewicz applies to this $(r+1)$-connected pair of simply
connected spaces and identifies the last group with
$\pi_{r+2}(Y,X)=\pi_{r+1}F$.  This proves (15), and the homotopy
exact sequence gives (16).  For normal crossings under (i), (8)
identifies $H_1$ of each mapping component with its meridian
kernel, proving the interpretation claimed.

Finally let
$Q^m=\{z_0^2+\cdots+z_m^2=1\}\subset\C^{m+1}$ and take
$(X,D)=(Q^m\times\C,Q^m\times\{0\})$.  Writing $z=x+iy$ gives
$x\cdot y=0$, $|x|^2=1+|y|^2$; hence
$(x,y)\mapsto(x/\sqrt{1+|y|^2},y)$ is a diffeomorphism
$Q^m\cong TS^m$.  Thus $Q^m\simeq S^m$.  For $m\geq2$ it is
simply connected.  The ordinary exit category is that of
$(\C,0)$: the closed object has trivial automorphisms, the open
object has group $\Z$, and there is a unique arrow from the
closed to the open object, fixed by the target action.  The closed
object is initial, so its classifying space is contractible.
The source has type $S^m$, and the first failed group is exactly
$\pi_m=\Z$.  The normal meridian is still injective, so these
examples specifically isolate failures of asphericity.  They also
show why one cannot replace (15)--(16) by a claim that every
failure of the criterion first changes $\pi_2$ of the ambient
space.
\end{proof}

\begin{proposition}[\status{Proved}: faces, tubes, and flags]\label{P1a2}
For manifolds, the real oriented blow-up is locally
\[
 M=X^{\log}=\{(r_i,\theta_i,w):r_i\geq0,\ \theta_i\in S^1\},
 \qquad z_i=r_i\theta_i.                                   \tag{3}
For stacks this formula is interpreted by étale descent.  Over $D_d^\circ$
its open codimension-$d$ face is canonically
$P_d^*/(\R_{>0})^d$.  The quotient is by the corresponding positive-real
subgroup of the possibly twisted torus, not by globally labelled factors.
The interior inclusion into each normalized closed face is a homotopy
equivalence.  For a flag, the inclusion of its deeper open face into the
shallower closed face, followed by the inverse of this equivalence, induces
Grothendieck's generization morphism on fundamental groupoids
(147/133--139, (219)--(222)).

The punctured tubular germ of a normalized branch intersection $D'$
has the homotopy type of its punctured multinormal bundle over the open
part $D'^\circ$:
\[
 \Pi_\infty\mathcal V^*_{D',X}\simeq
 \Pi_\infty P_{D'^\circ}.                                  \tag{4}
This implies (239), 147/145, including the stack case without a global
tubular embedding.  The equivalence is natural as a specialization
zigzag; base-point comparisons have the usual path and conjugacy choices.

Let $\mathcal A=\Exit(X,D)$, $\mathcal A_D$ its closed part, and
$\mathcal A_U=\Pi_\infty U$.  The crossing category is
\[
 \mathcal H=\operatorname{Fun}(\Delta^1,\mathcal A)
       \mathop{\times}_{\mathcal A\times\mathcal A}
       (\mathcal A_D\times\mathcal A_U).
\]
The space of $r$ composable morphisms in $\mathcal H$, with stratum flags
specified and equivalences in the endpoints retained, has homotopy type
$P^*_{d_0,\ldots,d_r}$.  Hence $\widetilde\Pi_{d_*}$ is the
\emph{degreewise fundamental groupoid} of this coherent nerve.  It is not
in general the nerve of the crossing category of $h\mathcal A$.
Finally,
\[
 L\simeq B\mathcal H\simeq
 \hocolim_{d_0>\cdots>d_r>0}P^*_{d_0,\ldots,d_r},           \tag{5}
\]
where the last expression is the semisimplicial realization with all
coherent face maps.  Inserting repeated ranks gives the equivalent full
simplicial nerve with degeneracies.  The internal flag category,
not $M$, is the ``déploiement'' of 147/114.
\end{proposition}
\begin{proof}
Formula (3) is coordinate independent: a change $z'_i=u_i z_{\sigma(i)}$
changes the boundary angle by $u_i/|u_i|$ and permutes the factors.  It is
also the divisorial instance of the definition of
\citet[\S1.2]{KatoNakayama1999}.  On an open face the nonzero normal
vectors have independent positive lengths, and dividing them out gives
the unit-torus bundle.  In a manifold with corners an inward pointing
vector field, multiplied by a sufficiently small positive function on a
locally finite chart cover, has a short flow defined everywhere.  Its
flow takes a closed face into its interior and is homotopic to the
identity.  The same flow restricted to the interior stays there.  This
proves the asserted homotopy equivalence, also for a noncompact face.

The punctured tube corresponds in $M$ to a collar of the normalized closed
face.  Projection along its positive normal radii is a homotopy
equivalence onto that face.  Composing with the preceding interior
equivalence proves (4) for a manifold.  The construction treats the
additional branches meeting $D'$ as boundary faces, which is why its
right side lies over $D'^\circ$ and not over the closed intersection.

Here is a descent formulation that does not assume compatible global
tubular embeddings.  Use the germ of deformation to the normal cone of
the local immersion.  In adapted coordinates its specialization chart
is $(s,v,t)\mapsto(s,tv,t)$, with the required $v_i$ nonzero.  At $t=0$
one obtains the punctured normal bundle and at $t>0$ the punctured tube
germ.  On every coordinate intersection both fibre inclusions into this
germ over the contractible interval are weak equivalences: contract the
positive radii in these coordinates.  This is a natural zigzag on an
étale hypercover; the deformation construction itself commutes with
étale pullback.  Realizing the hypercover preserves its termwise weak
equivalences and proves (4) on the stack.  Equivalently, a covering is
continued along a nonzero normal tangent ray; coverings and this
continuation descend uniquely.  This supplies the equivalence on
$\Pi_1$ without any chosen invariant splitting of a tangent sequence.
The transvection example in scan 146 therefore obstructs an embedding,
not this equivalence.

For the crossing assertion, a string of squares between arrows into
$U$ can be contracted along its target row, since that row lies in a
groupoid.  What remains is a string
$S_0\to S_1\to\cdots\to S_r\to U$.  Pull back the successive link
bundles of (1).  Their relative ranks add to $d_0$, and their positive
radial choices form a contractible space.  Their product is precisely
the multinormal torus bundle over the incidence cover
$D_{d_0,\ldots,d_r}^\circ$.  Dropping a middle term composes exits;
dropping the first uses the closed-face generization just described.
This proves the assertion including its coherences.

The normalized closed faces of $M$ cover $\partial M$.  Their multiple
intersections are indexed locally by nonempty subsets of the boundary
branches.  Barycentrically refining this incidence resolution replaces
such subsets by their strict chains and gives exactly the flag incidence
faces.  Their inclusions are
cofibrations in corner charts.  The augmented Čech realization is
therefore equivalent to their union: the fibre over a point is the
simplex on the finite set of local boundary faces containing it, hence
contractible; collars give the required local contraction.  Replacing
each face by its interior uses a termwise equivalence, proving (5).
The same Čech argument on a stack presentation gives stack homotopy
types, not coarse orbit spaces.  For an ordinary categorical crossing
category the simplicial-link identification in
\citet[Theorem 6.1]{DouteauWaas2022} applies after identifying its nerve
with $\operatorname{Map}_{\Delta^1}(\Delta^1,N\mathcal A)$.
That result justifies the STRATA C.2 comparison in that special setting;
it does not license replacing $\mathcal H$ by its level-zero truncation.
\end{proof}

\subsection{Meridians and the one-type criterion}

\begin{theorem}[\status{Proved}: the exact ordinary-exit criterion]\label{P1b1}
For every normal-crossings manifold or analytic DM stack in our scope,
the following are equivalent:
\begin{enumerate}[label=(\roman*)]
\item $\Exit(X,D)$ is equivalent to the nerve of an ordinary category.
\item Every connected stratum has a $K(\pi,1)$ homotopy type, and for
 every geometric boundary point $x$, every local branch choice and every
 path to a base point $u$ in the appropriate component of $U$, the map
 \[
  \mu_x:\Z^{I(x)}\longrightarrow\pi_1(U,u)                 \tag{6}
 \]
 induced by a small oriented normal torus is injective.
\item Every stratum is aspherical, and every relative torus lattice
 for every incidence $S<T$ injects into $\pi_1(T,t)$.
\end{enumerate}
Injectivity is independent of the indicated connecting paths.  Moreover,
the open-stratum condition (6) implies the relative condition without
assuming asphericity.  This proves the precise categorical part of
Lead 2 and 149/130; it does not assert that (6) is needed for every
possible derived reconstruction from the groupoid diagram.
\end{theorem}
\begin{proof}
An $\infty$-category is ordinary exactly when all its mapping spaces are
homotopy-discrete.  On one stratum the endomorphism space is $\Omega S$;
it is homotopy-discrete exactly when $\pi_j(S)=0$ for $j\geq2$.
Suppose henceforth that the strata are aspherical.  Each component $F$
of $F_{ST}$ is also aspherical, because it is a covering of $S$.
The fibration $(S^1)^r\to L_F\to F$, $r=d-e$, has exact sequence
\[
 1\longrightarrow\Z^r\longrightarrow G_F:=\pi_1(L_F)
   \longrightarrow\pi_1(F)\longrightarrow1,               \tag{7}
\]
and $L_F$ is aspherical, by its homotopy exact sequence.  The covering
map embeds $\pi_1(F)$ in $\pi_1(S)$.  By (2), each nonempty mapping-space
component is
\[
 K\!\left(\ker\{G_F\longrightarrow
       \pi_1(S)\times\pi_1(T)\},1\right).                 \tag{8}
Indeed the exact homotopy sequence of the homotopy fibre gives this
kernel as $\pi_1$ and makes all higher homotopy groups zero.  The kernel
of the first projection of the homomorphism in (8) is exactly the
fibre group $\Z^r$.  Therefore the kernel in (8) is the kernel of the
relative meridian homomorphism $\Z^r\to\pi_1(T)$.  This proves the
equivalence of the first and third conditions.  Applying it to $T=U$
proves that the first implies the second.

For the remaining implication we prove the stronger assertion stated
last.  At $x\in S$ let $I$ be its full branch set and let $J\subset I$
be the branches containing the selected germ of $T$.  Put
$L_T=P_T/(\R_{>0})^{|J|}$, the normal torus bundle over $T$.  In one
adapted coordinate chart there is a homomorphism
\[
 \Z^I\xrightarrow{a}\pi_1(L_T)\xrightarrow{b}\pi_1(U),
 \qquad ba=\mu_x,                                        \tag{9}
\]
obtained by displacing the coordinates in $I\setminus J$.  Its composite
to $\pi_1(T)$ is the relative meridian map on
$\Z^{I\setminus J}$ and zero on $\Z^J$.  The homotopy exact sequence
of $L_T\to T$ says
\[
 \ker\{\pi_1(L_T)\to\pi_1(T)\}
       =\operatorname{im}\{\Z^J\to\pi_1(L_T)\}.            \tag{10}
If $v\in\Z^{I\setminus J}$ maps to $1$ in $\pi_1(T)$, (10) gives
$w\in\Z^J$ with $a(v)=a(w)$.  Thus $\mu_x(v-w)=1$.
Injectivity of (6) implies $v-w=0$ in the \emph{domain} decomposition
$\Z^I=\Z^{I\setminus J}\oplus\Z^J$, so $v=w=0$.
This argument does not assume that the image of a torus lattice is a
direct summand of any fundamental group; that assertion in the lead
would be unjustified.  Coordinate paths specify (9); changing them
simultaneously conjugates the homomorphisms and leaves their kernels
unchanged.  The argument applies after any incidence cover and to the
stack homotopy fibre sequence, completing all cases.
\end{proof}

\begin{proposition}[\status{Refuted}: adjacent incidences are insufficient]\label{P1b2}
Even if every stratum is aspherical, injectivity for every incidence of
relative complex codimension one does not imply (6).  This already
fails for a smooth quasi-projective complex surface with a simple
normal-crossing divisor.
\end{proposition}
\begin{proof}
Start with the two axes in $\C^2$.  Blow up $p=(0,1)$ on the branch
$A=\{x=0\}$, denote the exceptional curve by $E$, and let
$q=E\cap A'$.  Delete one point $r\in E\setminus\{q\}$ from the
ambient blown-up surface.  Take $D=A'\cup B'\cup(E\setminus\{r\})$.
The open complement is unchanged: $U=(\C^*)^2$ with group
$\Z e_x\oplus\Z e_y$.  The other strata are
\[
 A'^\circ\simeq\C\setminus\{0,1\},\qquad
 B'^\circ\simeq\C^*,\qquad E^\circ\simeq\C^*,
\]
and the two point strata $o=A'\cap B'$ and $q=A'\cap E$.
Thus every stratum is aspherical.  The three rank-one meridians into
$U$ are $e_x,e_y,e_x$, respectively, so all are injective.
Every point-to-curve meridian is a puncture loop in a curve with at
least two punctures in its compactification; it has infinite order.
These are all relative-codimension-one incidences.

At $q$, however, both local boundary meridians map to $e_x$.  In a
blow-up chart $x=uv$, $y-1=v$, the two equations $u=0$ and $v=0$
exhibit this directly.  The lattice map is
$(a,b)\mapsto(a+b)e_x$ and has kernel $\Z(1,-1)$.
At the other double point it is the identity matrix on $\Z^2$.
The integer matrices and the punctured-curve ranks are checked in
\texttt{code/criterion\_checks.py}.  The deleted point is essential to
this example: before deleting it $E^\circ=\C$, and its adjacent test
already fails.  No properness of $X$ was assumed in the question.
\end{proof}

\begin{theorem}[\status{Proved}: the left-hand term and the meaning of $D_J$]\label{P1b3}
For a connected component $S$ of an open flag stratum and a rank-$r$
punctured normal torus bundle $P\to S$, based compatibly, the exact
sequence is
\[
 \pi_2(S)\xrightarrow{\partial}\Z^r\xrightarrow{\iota}
 \pi_1(P)\longrightarrow\pi_1(S)\longrightarrow1.           \tag{11}
\]
Consequently its kernel group is $\Z^r/\operatorname{im}\partial$,
which can have torsion and need not be a lattice,
and a zero can be inserted on the left exactly when $\partial=0$.
After pulling back to a sphere, the $i$th coordinate of $\partial$ is
the first Chern number of the $i$th normal line.  The assertion is valid
for stack homotopy groups and twisted branch lattices as well.

The topological form of 149/(37) for the trace of $U$ on a tube has
base $D_J^\circ$, the open intersection in (36).  It is not, in general,
exact with base the closed intersection $\overline D_J$ and kernel
generated only by the branches in $J$.  If its lattice map to
$\pi_1(U)$ is injective, then $\partial=0$, exactly as argued in
149/134.
\end{theorem}
\begin{proof}
The fibre $(\C^*)^r$ is connected with $\pi_1=\Z^r$ and no higher
homotopy.  The homotopy exact sequence is (11), with the usual
$\pi_1(S)$-action on the lattice.  Pulling back along a based sphere
trivializes the branch permutation local system.  The clutching map
on its equator has $i$th winding number equal to the Chern number;
the definition of the boundary homomorphism by lifting the two
hemispheres gives precisely that tuple.  This proves both the
necessity and sufficiency of $\partial=0$, without replacing it by
the stronger hypothesis $\pi_2(S)=0$.

For example, the normal circle bundle of a projective line in
$\PP^2$ is the Hopf bundle.  Here $\partial:\Z\to\Z$ is multiplication
by $1$, and $\pi_1(P)=1$.  For the trivial normal line of
$\PP^1\times\{0\}\subset\PP^1\times\C$, $\pi_2(S)=\Z$ but
$\partial=0$ and the meridian injection is exact on the left.

To see the closed-intersection error without any $\pi_2$ complication,
use the axes in $\C^2$ and $J=\{x=0\}$.  A punctured tube traced on
$U=(\C^*)^2$ has group $\Z^2$.  Its open base is $\C^*$ and its
sequence is $1\to\Z\to\Z^2\to\Z\to1$.  If the base is changed to
the closed axis $\C$, the claimed kernel would be $\Z$, whereas the
actual kernel of $\Z^2\to1$ has rank two.  The missing generator
goes around the other branch on the base.  A logarithmic closed base
could retain that generator, but an ordinary closed base cannot.
Finally, $\iota\partial=0$ in (11); its composite with
$\pi_1(P)\to\pi_1(U)$ is zero.  Injectivity of the latter composite
on $\Z^r$ forces $\partial=0$.  This is the exact content of the
manuscript's argument, independent of either doubtful decoration on
$D_J$.  The Hopf and rank computations are included in the program.
\end{proof}
